\chapter{هندسه تحلیلی}

می‌دانیم از هر دو نقطۀ متمایز، تنها یک خط عبور می‌کند؛ بنابراین:
\begin{enumerate}
    \item با داشتن \underline{مختصات دو نقطه} از یک خط باید بتوان معادلۀ آن را بدست آورد.
    \item با داشتن معادلۀ یک خط می‌توان با مشخص کردن \underline{دو نقطه} از خط، نمودار آن را در دستگاه محورهای مختصات رسم نمود.
\end{enumerate}

\begin{example}
نمودار خطوط با معادلات زیر را در دستگاه محورهای مختصات مقابل رسم کنید.
\end{example}

\begin{LTR}
\begin{minipage}[c]{0.35\textwidth}
\vspace{0pt}
\begin{enumerate}
\setlength{\itemsep}{0.7cm}
\item $L_1:y=2x-4$
\item $L_2:y=2x+1$
\item $L_3:y=1$
\item $L_4:x=-2$
\end{enumerate}
\end{minipage}\hfill\begin{minipage}[c]{0.55\textwidth}
    \vspace{0pt}
    \centering

    \begin{tikzpicture}[
        x=0.7cm,
        y=0.7cm,
        line cap=round,
        >=stealth
    ]

        % Grid
        \draw[
            step=1,
            color=MPBorder,
            line width=0.35pt
        ] (-3,-2) grid (3,3);

        % X axis
        \draw[
            ->,
            color=MPPrimary,
            line width=0.9pt
        ] (-4.6,0) -- (5.6,0)
        node[below right] {$x$};

        % Y axis
        \draw[
            ->,
            color=MPPrimary,
            line width=0.9pt
        ] (0,-2.4) -- (0,3.5)
        node[above left] {$y$};

        % X-axis labels
        \foreach \x in {-3,-2,-1,1,2,3} {
            \node[
                below=2pt,
                fill=MPBackground,
                inner sep=1pt,
                text=MPTextMuted
            ] at (\x,0) {\small $\x$};
        }

        % Y-axis labels
        \foreach \y in {-2,-1,1,2,3} {
            \node[
                left=2pt,
                fill=MPBackground,
                inner sep=1pt,
                text=MPTextMuted
            ] at (0,\y) {\small $\y$};
        }

    \end{tikzpicture}
\end{minipage}
\end{LTR}

\vspace{1.8\baselineskip}

% شیب خط
\begin{minipage}[c]{0.46\textwidth}

    \section{شیب (ضریب زاویه) خط:}

\end{minipage}\hfill\begin{minipage}[c]{0.22\textwidth}
\centering
\[
    \text{شیب}=\frac{a}{b}
\]
\end{minipage}\hfill\begin{minipage}[c]{0.24\textwidth}
\centering
\begin{tikzpicture}[
    x=0.65cm,
    y=0.65cm,
    line cap=round
]

    % Triangle vertices
    \coordinate (A) at (0,0);
    \coordinate (B) at (3,0);
    \coordinate (C) at (3,1.8);

    % Main line
    \draw[
        MPPrimary,
        line width=1.3pt
    ] (-0.3,-0.18) -- (3.3,1.98);

    % Horizontal dashed line: b
    \draw[
        MPSecondary,
        dashed,
        line width=0.9pt
    ] (A) -- (B);

    % Vertical dashed line: a
    \draw[
        MPSecondary,
        dashed,
        line width=0.9pt
    ] (B) -- (C);

    % Right-angle marker
    \draw[
        MPText,
        line width=0.7pt
    ]
    (2.72,0) -- (2.72,0.28) -- (3,0.28);

    % Labels
    \node[
        below=3pt,
        text=MPSecondary
    ] at (1.5,0) {$b$};

    \node[
        right=3pt,
        text=MPSecondary
    ] at (3,0.9) {$a$};

\end{tikzpicture}
\end{minipage}

\begin{enumerate}[label=\arabic*]
\item شیب یک خط برابر با تانژانت زاویه‌ای است که خط با جهت مثبت محور
$x$ها می‌سازد که آن را با $m$ نشان می‌دهیم.

% Positive slope: m = tan(alpha)
\begin{minipage}[c]{0.19\textwidth}
    \centering
    $\displaystyle m=\tan\alpha$

    \vspace{0.4em}

    $\displaystyle m>0$
\end{minipage}\hfill\begin{minipage}[c]{0.26\textwidth}
    \centering
    \begin{tikzpicture}[
        x=0.65cm,
        y=0.65cm,
        >=stealth
    ]
        % Axes
        \draw[->, MPPrimary, thick]
            (-4,0) -- (2,0);

        \draw[->, MPPrimary, thick]
            (-1,-1.7) -- (-1,1.7);

        % Positive-slope line
        \draw[MPPrimary, line width=1.5pt]
            (-0.9,-1.4) -- (1.1,1.6);

        % Angle alpha
        \draw[MPSecondary, thick]
            (0.033,0) ++(0:0.35)
            arc[start angle=0,end angle=56,radius=0.35];

        \node[MPPrimary] at (0.65,0.42) {$\alpha$};
    \end{tikzpicture}
\end{minipage}\hfill\begin{minipage}[c]{0.26\textwidth}
    \centering
    \begin{tikzpicture}[
        x=0.65cm,
        y=0.65cm,
        >=stealth
    ]
        % Axes
        \draw[->, MPPrimary, thick]
            (-4,0) -- (2,0);

        \draw[->, MPPrimary, thick]
            (-1,-1.7) -- (-1,1.7);

        % Negative-slope line
        \draw[MPPrimary, line width=1.5pt]
            (-0.9,1.4) -- (1.1,-1.6);

        % Angle beta
        \draw[MPSecondary, thick]
            (0.033,0) ++(0:0.35)
            arc[start angle=0,end angle=124,radius=0.35];

        \node[MPPrimary] at (0.15,0.65) {$\beta$};
    \end{tikzpicture}
\end{minipage}\hfill\begin{minipage}[c]{0.19\textwidth}
    \centering
    $\displaystyle m=\tan\beta$

    \vspace{0.4em}

    $\displaystyle m<0$
\end{minipage}

\begin{minipage}[c]{0.7\textwidth}
\item شیب خطی که از دو نقطۀ
$A\left|\begin{matrix}x_1\\y_1\end{matrix}\right.$
و
$B\left|\begin{matrix}x_2\\y_2\end{matrix}\right.$
می‌گذرد برابر است با:
\end{minipage}\hfill\begin{minipage}[c]{0.15\textwidth}
    \[
        m=\frac{y_2-y_1}{x_2-x_1}
    \]
\end{minipage}

\item اگر شیب یک خط \(m=0\) باشد، خط موازی محور \(x\)ها است و معادلۀ آن به‌صورت \(y=a\) است.

\begin{center}
\begin{tikzpicture}[
    x=0.75cm,
    y=0.75cm,
    >=stealth
]
    % Coordinate axes
    \draw[->, MPPrimary, line width=0.9pt]
        (-2.5,0) -- (2.5,0);

    \draw[->, MPPrimary, line width=0.9pt]
        (0,-0.5) -- (0,1.5);

    % Horizontal line
    \draw[MPPrimary, line width=1.1pt]
        (-1.8,0.85) -- (1.8,0.85);

    % Equation label
    \node[right] at (1.8,0.85) {$y=a$};

\end{tikzpicture}
\end{center}

\item اگر شیب یک خط تعریف‌نشده باشد، خط موازی محور \(y\)ها است و معادلۀ آن به‌صورت \(x=k\) است.

\begin{center}
\begin{tikzpicture}[
    x=0.75cm,
    y=0.75cm,
    >=stealth
]
    % Coordinate axes
    \draw[->, MPPrimary, line width=0.9pt]
        (-2.5,0) -- (2.5,0);

    \draw[->, MPPrimary, line width=0.9pt]
        (0,-1.5) -- (0,1.5);

    % Vertical line
    \draw[MPPrimary, line width=1.1pt]
        (0.9,-1.5) -- (0.9,1.5);

    % Equation label
    \node[below right] at (0.9,-1.5) {$x=k$};

\end{tikzpicture}
\end{center}

\item \begin{enumerate}[label=\harfi*]
    \item اگر معادلۀ خط به‌صورت \(ax+by+c=0\) (\(x\) و \(y\) یک طرف مساوی) باشد، آنگاه:
    \hfill
    \begin{minipage}[c]{0.154\textwidth}
    \centering
    $$m=-\frac{x \text{ضریب}}{y \text{ضریب}}$$
    \end{minipage}

    \item اگر معادلۀ خط به‌صورت \(ax=by+c\) (\(x\) و \(y\) در دو طرف مساوی) باشد، آنگاه:
    \hfill
    \begin{minipage}[c]{0.14\textwidth}
    \centering
    $$m=\frac{x \text{ضریب}}{y \text{ضریب}}$$
    \end{minipage}
\end{enumerate}

\vspace{0.4\baselineskip}

\begin{minipage}[c]{0.6\textwidth}
\item دو خط \(d_1\) و \(d_2\) وقتی با هم موازی‌اند که:
\hfil $m_1=m_2$
\end{minipage}
\hfill
\begin{minipage}[c]{0.3\textwidth}
    \centering

    \begin{tikzpicture}[
        x=0.65cm,
        y=0.65cm,
        >=stealth
    ]

        % First parallel line
        \draw[->, MPPrimary, line width=1.1pt] (-2,0.4) -- (1,1.3);
        \draw[MPPrimary, line width=1.1pt] (1,1.3) -- (2.3,1.7);

        % Second parallel line
        \draw[->, MPPrimary, line width=1.1pt] (-2,-0.6) -- (1,0.3);
        \draw[MPPrimary, line width=1.1pt] (1,0.3) -- (2.3,0.7);
        
        % Line labels
        \node[left] at (-2,0.4) {$d_1$};
        \node[left] at (-2,-0.6) {$d_2$};

        % Slope labels
        \node[above] at (0,1) {$m_1$};
        \node[above] at (0,0) {$m_2$};

    \end{tikzpicture}
\end{minipage}

% ==========

\begin{minipage}[c]{0.4\textwidth}
\item دو خط \(d_1\) و \(d_2\) وقتی بر هم عمودند که:
\end{minipage}
\hfill
\begin{minipage}[c]{0.15\textwidth}
    \centering
    \[
        m_1m_2=-1
    \]
\end{minipage}
\hfill
\begin{minipage}[c]{0.03\textwidth}
    \centering
    یا
\end{minipage}
\hfill
\begin{minipage}[c]{0.32\textwidth}
    \centering
    \[
        m_1=-\frac{1}{m_2}
    \]

    {\small
        (شیب‌ها عکس و قرینۀ یکدیگرند)
    }
\end{minipage}

\begin{center}
\begin{tikzpicture}[
    x=0.65cm,
    y=0.65cm
]

    % First line
    \draw[
        MPPrimary,
        line width=1.1pt
    ] (-1.4,-1.4) -- (1.4,1.4);

    % Second line
    \draw[
        MPPrimary,
        line width=1.1pt
    ] (-1.4,1.4) -- (1.4,-1.4);

    % Right angle marker
    \draw[
        MPPrimary,
        line width=0.8pt
    ]
    (0.17,0.17) --
    (0.34,0) --
    (0.17,-0.17);

    % Line labels
    \node[left] at (-1.4,-1.4) {$d_1$};
    \node[left] at (-1.4,1.4) {$d_2$};

    % Slope labels
    \node[below] at (-0.45,-0.45) {$m_1$};
    \node[above] at (-0.45,0.45) {$m_2$};

\end{tikzpicture}
\end{center}

\item شرط آنکه سه نقطۀ \(A\)، \(B\) و \(C\) بر یک خط راست واقع باشند (بر یک استقامت یا در یک راستا باشند)، آن است که:

\begin{minipage}[c]{0.19\textwidth}
    \centering
    \[
        m_{AB}=m_{BC}
    \]
\end{minipage}
\hfill
\begin{minipage}[c]{0.33\textwidth}
    \centering

    \begin{tikzpicture}[
        x=0.65cm,
        y=0.65cm
    ]

        % Collinear dotted line
        \draw[
            MPPrimary,
            line width=1pt,
            dotted
        ] (-2,0) -- (2.2,0.9);

        % Points
        \fill[MPPrimary] (-1.3,0.16) circle (1.5pt);
        \fill[MPPrimary] (0,0.45) circle (1.5pt);
        \fill[MPPrimary] (1.3,0.75) circle (1.5pt);

        % Point labels
        \node[above] at (-1.3,0.16) {$A$};
        \node[above] at (0,0.45) {$B$};
        \node[above] at (1.3,0.75) {$C$};

    \end{tikzpicture}
\end{minipage}
\end{enumerate}

\begin{example}
زاویۀ بین خط
$y=\sqrt{3}x-1$
و جهت مثبت محور $x$ها کدام است؟

\begin{tasks}(4)
    \task $120^\circ$
    \task $60^\circ$
    \task $45^\circ$
    \task $30^\circ$
\end{tasks}
\end{example}

\begin{remark}[یادآوری]
\hfill
$\tan 30^\circ = \frac{\sqrt{3}}{3} \qquad \tan 45^\circ = 1 \qquad \tan 60^\circ = \sqrt{3}$
\end{remark}


\begin{example}
به ازای چه مقداری از $m$ دو خط $y=2x+1$ و $y+mx+2=0$ بر هم عموداند؟

\begin{tasks}(4)
    \task $-2$
    \task $-\frac12$
    \task $\frac12$
    \task $2$
\end{tasks}
\end{example}

\begin{example}
اگر سه نقطۀ متمایز $A(1,1)$، $B(-1,-1)$ و $C(m,m^2+2m)$ روی یک خط راست باشند آن‌گاه:

\begin{tasks}(4)
    \task $m=0$
    \task $m=-1$
    \task $m=1$
    \task $m=0\text{ و }-1$
\end{tasks}
\end{example}

\vspace{0.4\baselineskip}

\begin{remark}
مختصات نقطۀ $M$ وسط $AB$ برابر است با:
\hfill\begin{minipage}[c]{0.2\textwidth}
$$M\left|\begin{matrix}
\frac{x_{\! \scriptscriptstyle A}+x_{\! \scriptscriptstyle B}}{2}
\\
\frac{y_{\! \scriptscriptstyle A}+y_{\! \scriptscriptstyle B}}{2}
\end{matrix}\right.$$
\end{minipage}
\end{remark}

\begin{remark}
طول نقطۀ برخورد خط با محور $x$ها (طول‌ها) را \underline{طول از مبدأ} و عرض نقطۀ برخورد خط با محور $y$ها (عرض‌ها) را \underline{عرض از مبدأ} خط می‌گویند.
\end{remark}

\begin{minipage}[c]{0.55\textwidth}
برای بدست آوردن طول از مبدأ خط کافی است در معادلۀ خط بجای $y$، صفر قرار دهیم و برای تعیین عرض از مبدأ در معادلۀ خط بجای $x$، صفر قرار می‌دهیم.

\[
\text{\rl{معادلۀ خط d}}:\frac{x}{p}+\frac{y}{q}=1
\]
\end{minipage}
\hfill
\begin{minipage}[c]{0.35\textwidth}
\centering
\begin{tikzpicture}[
        x=0.6cm,
        y=0.4cm,
        >=stealth
    ]

    % fill
    \fill[fill=MPSecondary!20] (0,0)--(0,2)--(-2,0);

    % Axes
    \draw[->, MPPrimary, thick]
        (-3.2,0) -- (3.2,0);

    \draw[->, MPPrimary, thick]
        (0,-1.8) -- (0,3.2);

    % Positive-slope line
    \draw[MPPrimary, line width=1.5pt]
        (-3,-1) -- (1,3);

    \node[above left] at (0,2) {$q$};
    \node[below right] at (-2,0) {$p$};
    \node[left] at (-3,-1) {$d$};
\end{tikzpicture}
    \[
    S=\frac{|p\times q|}{2}
    \]
\end{minipage}

\begin{remark}
معادلۀ خطی که از نقطۀ
$M\left|\begin{matrix}
x_1
\\
y_1
\end{matrix}\right.$
گذشته و دارای شیب $m$ باشد بصورت زیر است:
\[y-y_1=m(x-x_1)\]
\end{remark}

\begin{example}
نقطۀ $m(5,-4)$ وسط پاره‌خط واصل بین دو نقطۀ $A$ و $B(7,-2)$ است. مختصات $A$ را بیابید.
\end{example}

\vspace{0.8\baselineskip}

\begin{example}
معادلۀ خطی که از نقطۀ $A(-1,-5)$ گذشته و طول از مبدأ آن 4 باشد، کدام است؟

\begin{tasks}(2)
    \task $x+y=4$
    \task $x-y=4$
    \task $y-x=4$
    \task $x+y=-4$
\end{tasks}
\end{example}

\begin{example}
دو نقطه‌ی $A(14,3)$ و $B(10,-13)$ مفروض است، معادله‌ی عمود منصف پاره خط $AB$ را بنویسید.
\end{example}

\vspace{0.8\baselineskip}

\begin{example}
خطی محورهای مختصات را در نقاط $A$ و $B$ قطع می‌کند.
اگر $M(-3,2)$ وسط $AB$ باشد. معادلۀ این خط کدام است؟

\begin{tasks}(2)
    \task $2x-3y=-12$
    \task $2x-3y=12$
    \task $3x+2y=12$
    \task $3x-2y=-12$
\end{tasks}
\end{example}

\begin{example}
خطی که از نقاط \((-1,1)\) و \((2,2)\) می‌گذرد،
با محورهای مختصات چه مساحتی می‌سازد؟

\begin{tasks}(1)
    \task \(\frac{2}{3}\)
    \task \(\frac{16}{3}\)
    \task \(\frac{4}{3}\)
    \task \(\frac{8}{3}\)
\end{tasks}
\end{example}

\begin{remark}
\begin{minipage}[c]{0.34\textwidth}
    \centering
    \resizebox{0.92\linewidth}{!}{%
        \begin{tikzpicture}[>=stealth]
            % Coordinate axes
            \draw[->,MPPrimary,thick]
                (-1.5,0) -- (1.5,0);
            \draw[->,MPPrimary,thick]
                (0,-1.5) -- (0,1.5);

            % Positive slope
            \draw[
                MPPrimary,
                thick,
                dotted
            ] (-1.25,-1.25) -- (1.25,1.25);

            % Labels
            \node[rotate=45] at (0.85,0.95) {\tiny$y=x$};
            \node[rotate=45] at (-0.9,-0.65)
                {\tiny\rl{نیمساز ناحیه اول سوم}};
        \end{tikzpicture}%
    }
\end{minipage}
\hfill
\begin{minipage}[c]{0.34\textwidth}
    \centering
    \resizebox{0.92\linewidth}{!}{%
        \begin{tikzpicture}[>=stealth]
            % Coordinate axes
            \draw[->,MPPrimary,thick]
                (-1.5,0) -- (1.5,0);
            \draw[->,MPPrimary,thick]
                (0,-1.5) -- (0,1.5);

            % Negative slope
            \draw[
                MPPrimary,
                thick,
                dotted
            ] (-1.25,1.25) -- (1.25,-1.25);

            % Labels
            \node[rotate=-45] at (0.8,-1.0) {\tiny$y=-x$};
            \node[rotate=-45] at (-0.8,1.0)
                {\tiny\rl{نیمساز ناحیه دوم چهارم}};
        \end{tikzpicture}%
    }
\end{minipage}
\end{remark}

\begin{remark}[توجه]
مختصات نقطۀ برخورد دو خط با حل معادله‌های آن دو در یک دستگاه
معادلات به دست می‌آید.
\end{remark}

\begin{example}
محل برخورد دو خط به معادله‌های
\(y=x+2\) و \(my=x+n\)
روی محور \(ox\) قرار دارد. \(n\) کدام است؟

\begin{tasks}(2)
    \task \(-2\)
    \task \(-1\)
    \task \(1\)
    \task \(2\)
\end{tasks}
\end{example}

\begin{example}
معادلۀ خطی که از نقطۀ تلاقی دو خط
\(y=x+2\) و \(y=2x+3\)
گذشته و بر نیمساز ربع اول و سوم عمود است،
کدام می‌باشد؟

\begin{tasks}(2)
    \task \(y=-x+1\)
    \task \(y=-x\)
    \task \(y=x+2\)
    \task \(y=x-1\)
\end{tasks}
\end{example}

\begin{example}
نقطۀ \(A(7,6)\) رأس یک متوازی‌الاضلاع است که
دو ضلع آن منطبق بر دو خط به معادلات
\(2y-3x=11\) و \(3y+4x=8\)
می‌باشند. مختصات وسط قطر آن کدام است؟

\begin{tasks}(2)
    \task \((1,5)\)
    \task \((3,4)\)
    \task \((3,5)\)
    \task \((4,3)\)
\end{tasks}
\end{example}

% ==========================================
% EXAMPLE 13
% ==========================================

\begin{example}
اگر خطوط به معادلات
\(
    x+y=2,\quad 2x-y=1,\quad x-my=-1
\)
از یک نقطه بگذرند، \(m\) کدام است؟

\begin{flushright}
\begin{minipage}{0.38\textwidth}
    \begin{tasks}(2)
        \task \(-2\)
        \task \(-1\)
        \task \(1\)
        \task \(2\)
    \end{tasks}
\end{minipage}
\end{flushright}
\end{example}

% ==========================================
% PARALLELOGRAM PROPERTY
% ==========================================

\begin{remark}[توجه]
در هر متوازی‌الاضلاع، قطرها منصف یکدیگرند و 
در متوازی‌الاضلاع \(ABCD\) داریم:

\begin{minipage}[c]{0.48\textwidth}
    \[
    \left\{
    \begin{aligned}
        x_{_A}+x_{_C} &= x_{_B}+x_{_D}\\
        y_{_A}+y_{_C} &= y_{_B}+y_{_D}
    \end{aligned}
    \right.
    \]
\end{minipage}
\hfill
\begin{minipage}[c]{0.42\textwidth}
    \centering
    \resizebox{0.95\linewidth}{!}{%
        \begin{tikzpicture}[
            line cap=round,
            line join=round
        ]
            % Vertices
            \coordinate (A) at (0,2);
            \coordinate (B) at (3,2);
            \coordinate (C) at (4.4,0.5);
            \coordinate (D) at (1.4,0.5);

            % Parallelogram
            \draw[MPPrimary, line width=1.1pt]
                (A)--(B)--(C)--(D)--cycle;

            % Diagonals
            \draw[MPSecondary, dotted, thick]
                (A)--(C);
            \draw[MPSecondary, dotted, thick]
                (B)--(D);

            % Intersection
            \fill[MPPrimary] (intersection cs:
                first line={(A)--(C)},
                second line={(B)--(D)}
            ) circle (2.2pt);

            % Vertex labels
            \node[above left] at (A) {$A$};
            \node[above right] at (B) {$B$};
            \node[below right] at (C) {$C$};
            \node[below left] at (D) {$D$};
        \end{tikzpicture}%
    }
\end{minipage}
\end{remark}

% ==========================================
% EXAMPLE 14
% ==========================================

\vspace{0.6\baselineskip}

\begin{example}
نقاط \(A(2,3)\)، \(B(-1,0)\) و \(C(1,-2)\)
سه رأس از یک مستطیل \(ABCD\) هستند.
مختصات رأس چهارم آن را بیابید.
\end{example}

% ==========================================
% DISTANCE BETWEEN TWO POINTS
% ==========================================

\begin{remark}[توجه]
فاصلۀ عددی دو نقطۀ
\(A(x_1,y_1)\) و \(B(x_2,y_2)\)،
یعنی طول پاره‌خط \(AB\)، برابر است با:

\begin{center}
    \fbox{%
        \(\displaystyle
        AB=\sqrt{(x_2-x_1)^2+(y_2-y_1)^2}
        \)
    }
\end{center}

\begin{minipage}[c]{0.48\textwidth}
    اگر \(A\) و \(B\) دو نقطۀ هم‌عرض باشند:
    \[
        AB=|x_{_A}-x_{_B}|
    \]
\end{minipage}
\hfill
\begin{minipage}[c]{0.43\textwidth}
    اگر \(A\) و \(B\) دو نقطۀ هم‌طول باشند:
    \[
        AB=|y_{_A}-y_{_B}|
    \]
\end{minipage}

\vspace{0.4\baselineskip}

\begin{minipage}[c]{0.48\textwidth}
    \[
    \begin{gathered}
        \Point{A}{-1}{3},\, \Point{B}{5}{3}\,
        \Longrightarrow AB= \quad
    \end{gathered}
    \]
\end{minipage}
\hfill
\begin{minipage}[c]{0.43\textwidth}
    \[
    \begin{gathered}
        \Point{A}{1}{2},\, \Point{B}{1}{-3}\,
        \Longrightarrow AB= \quad
    \end{gathered}
    \]
\end{minipage}
\end{remark}

% ==========================================
% EXAMPLE 15
% ==========================================

\vspace{0.6\baselineskip}

\begin{example}
نقاط \(A(2,0)\)، \(B(5,4)\) و \(C(-2,3)\)
را در نظر بگیرید.

\begin{enumerate}
    \setlength{\itemsep}{0.8\baselineskip}

    \item محیط مثلث \(ABC\) را بیابید.

    \item به دو روش نشان دهید
    \(\overset{\triangle}{ABC}\)
    یک مثلث قائم‌الزاویه است.

    \item مساحت مثلث \(ABC\) چقدر است؟
\end{enumerate}
\end{example}

% ==========================================
% EXAMPLE 16
% ==========================================

\vspace{0.6\baselineskip}

\begin{example}
مثلث با رئوس
\(A(1,9)\)، \(B(3,1)\) و \(C(7,11)\)
را در نظر بگیرید. طول میانه \(AM\) و معادلۀ آن
را به دست آورید.
\end{example}

% ==========================================
% EXAMPLE 17
% ==========================================

\vspace{0.8\baselineskip}

\begin{example}
دایره‌ای به مرکز مبدأ مختصات، از نقطۀ
\(N(-6,8)\) می‌گذرد. قطر دایره را محاسبه کنید.
\end{example}

\section{فاصلۀ نقطه از خط}


% ==========================================
% DISTANCE FROM A POINT TO A LINE
% ==========================================

\noindent
\begin{minipage}[c]{0.63\textwidth}
    \raggedleft

    فاصلۀ نقطۀ \(\Point{A}{x_1}{y_1}\) از خط
    \(d: ax+by+C=0\)، برابر با طول \(AH\)،
    یعنی طول پاره‌خط عمودی است که از نقطۀ \(A\)
    بر خط رسم می‌شود. که این طول برابر است با:

    \[
        AH=
        \frac{|ax_1+by_1+C|}
        {\sqrt{a^2+b^2}}
    \]

    \hfill {\small (اثبات صفحه ۲۷)}
\end{minipage}
\hfill
\begin{minipage}[c]{0.30\textwidth}
    \centering
    \begin{tikzpicture}[
        x=0.85cm,
        y=0.85cm,
        line cap=round,
        line join=round
    ]

        % Line d
        \draw[
            MPPrimary,
            line width=1pt
        ] (-1.7,-0.425) -- (1.7,0.425);

        % Point A and perpendicular AH
        \coordinate (H) at (0,0);
        \coordinate (A) at (-0.25,1);

        \draw[
            MPPrimary,
            line width=1pt
        ] (A) -- (H);

        % Point A
        \fill[MPPrimary] (A) circle (1.5pt);

        % Right-angle marker
        \draw[
            MPPrimary,
            line width=0.75pt
        ]
        ($(H)+(0.126,0.032)$)
        --
        ($(H)+(0.094,0.158)$)
        --
        ($(H)+(-0.032,0.126)$);

        % Labels
        \node[above left] at (A) {$A$};
        \node[below] at (H) {$H$};
        \node[below left] at (-1.55,-0.39) {$d$};

    \end{tikzpicture}
\end{minipage}

\begin{remark}[توجه]
    \begin{minipage}[c]{0.63\textwidth}
        \raggedleft
        فاصلۀ مبدأ مختصات از خط
        \(ax+by+C=0\) برابر است با:
    \end{minipage}
    \hfill
    \begin{minipage}[c]{0.25\textwidth}
        \centering
        \[
            \frac{|C|}{\sqrt{a^2+b^2}}
        \]
    \end{minipage}
\end{remark}

% ==========================================
% EXAMPLE 1
% ==========================================

\vspace{0.6\baselineskip}

\begin{example}
    فاصلۀ نقطۀ \(A(7,5)\) را از خط
    \(4x+3y-18=0\) به دست آورید.
\end{example}

% ==========================================
% EXAMPLE 2
% ==========================================

\vspace{0.8\baselineskip}

\begin{example}
    اگر \(A(-1,2)\)، \(B(3,0)\) و \(C(1,-2)\)
    سه رأس مثلث \(ABC\) باشند،
    طول ارتفاع \(AH\) کدام است؟

    \begin{tasks}(2)
        \task \(2\sqrt{3}\)
        \task \(3\sqrt{3}\)
        \task \(2\sqrt{2}\)
        \task \(3\sqrt{2}\)
    \end{tasks}
\end{example}

% ==========================================
% EXAMPLE 3
% ==========================================

\vspace{0.8\baselineskip}

\begin{example}
    مساحت دایره‌ای به مرکز \(W(2,-1)\)
    که بر خط \(3x-4y=0\) مماس است را بیابید.
\end{example}

% ==========================================
% EXAMPLE 4
% ==========================================

\vspace{0.8\baselineskip}

\begin{example}
    اگر فاصلۀ \(M(1,2)\) از خط
    \(3x+my-1=0\) برابر \(2\) باشد، آن‌گاه:

    \begin{tasks}(2)
        \task \(m=-2\)
        \task \(m=1\)
        \task \(m=4\)
        \task \(m=-1\)
    \end{tasks}
\end{example}


% ==========================================
% EXAMPLES: DISTANCE FROM A POINT TO A LINE
% ==========================================

% ==========================================
% EXAMPLE 5
% ==========================================

\vspace{0.6\baselineskip}

\begin{example}
    فاصلۀ مبدأ مختصات از خط
    \(
        2mx+(1-m^2)y=1+m^2
    \)
    کدام است؟ \((m\ne 1)\)

    \begin{tasks}(2)
        \task \(\dfrac{1}{m}\)
        \task \(m\)
        \task \(1\)
        \task \(1+m^2\)
    \end{tasks}
\end{example}

% ==========================================
% EXAMPLE 6
% ==========================================

\vspace{0.2\baselineskip}

\begin{example}
    مساحت مربعی که یک ضلع آن بر خط به معادلۀ
    \(y=x-1\) واقع و مختصات یک رأس آن
    \((2,3)\) باشد، کدام است؟

    \begin{tasks}(2)
        \task \(2\)
        \task \(3\)
        \task \(4\)
        \task \(9\)
    \end{tasks}
\end{example}

% ==========================================
% EXAMPLE 7
% ==========================================

\vspace{0.2\baselineskip}

\begin{example}
    دو نقطه بر خط به معادلۀ \(y=x-1\) قرار دارند
    که فاصلۀ این نقاط از خط به معادلۀ
    \(2x-3y=5\)، برابر \(\sqrt{13}\) است.
    طول این دو نقطه کدام است؟

    \begin{tasks}(2)
        \task $9$ و $-15$
        \task $11$ و $-15$
        \task $-11$ و $+15$
        \task $-9$ و $11$
    \end{tasks}
\end{example}

% ==========================================
% EXAMPLE 8
% ==========================================

\vspace{0.2\baselineskip}

\begin{example}
    دو ضلع یک مستطیل منطبق بر دو خط به معادلات
    \(2y+x=6\) و \(2x-y=7\) هستند و یکی از
    رأس‌های آن نقطۀ \(A(8,5)\) است.
    مساحت این مستطیل کدام است؟

    \begin{tasks}(4)
        \task \(7.2\)
        \task \(9.6\)
        \task \(11.4\)
        \task \(12.8\)
    \end{tasks}
\end{example}

% ==========================================
% DISTANCE BETWEEN PARALLEL LINES
% ==========================================

\begin{remark}[توجه]
    \noindent
    \begin{minipage}[c]{0.68\textwidth}
        \raggedleft
        فاصلۀ دو خط موازی
        \(
            ax+by+c=0
        \)
        و
        \(
            ax+by+c'=0
        \)
        برابر است با:
    \end{minipage}
    \hfill
    \begin{minipage}[c]{0.25\textwidth}
        \centering
        \[
            \frac{|c-c'|}{\sqrt{a^2+b^2}}
        \]
    \end{minipage}
\end{remark}

\begin{center}
\begin{tikzpicture}[
    line cap=round,
    line join=round
]

    % Two parallel lines
    \draw[
        MPPrimary,
        line width=1pt,
        -{Stealth}
    ] (0,1.1) -- (2.1,1.1);

    \draw[
        MPPrimary,
        line width=1pt,
        -{Stealth}
    ] (0,0) -- (2.1,0);

    % Perpendicular distance
    \draw[
        MPPrimary,
        dashed,
        line width=0.9pt
    ] (1.05,0) -- (1.05,1.1);

    % Right-angle marker
    \draw[
        MPPrimary,
        line width=0.8pt
    ] (1.05,0.95) -- (1.2,0.95) -- (1.2,1.1);

\end{tikzpicture}

\vspace{0.3\baselineskip}
{\small (اثبات صفحه ۲۷)}    
\end{center}

\begin{example}
    دو ضلع مربع منطبق بر دو خط به معادلات
    \(2x-2y=3\) و \(y=x+1\) هستند.
    مساحت این مربع کدام است؟

    \begin{tasks}(2)
        \task \(\dfrac{9}{8}\)
        \task \(\dfrac{9}{4}\)
        \task \(\dfrac{25}{8}\)
        \task \(\dfrac{25}{4}\)
    \end{tasks}
\end{example}


% ==========================================
% LOCUS OF POINTS EQUIDISTANT FROM
% TWO PARALLEL LINES
% ==========================================

\begin{remark}
    مکان هندسی نقاطی از صفحه که از دو خط موازی
    \(ax+by+c=0\) و \(ax+by+c'=0\)
    به یک فاصله باشند، خطی موازی آن‌هاست که از
    وسط فاصلۀ بین دو خط می‌گذرد و معادلۀ آن
    به صورت زیر است:
\end{remark}

\begin{center}
    \begin{tikzpicture}[
        line cap=round,
        line join=round
    ]

        % Upper parallel line
        \draw[
            MPPrimary,
            line width=1pt
        ] (-2.6,1.15) -- (0.8,1.15);

        \node[
            above=2pt
        ] at (-0.9,1.15)
        {$ax+by+c=0$};

        % Lower parallel line
        \draw[
            MPPrimary,
            line width=1pt
        ] (-2.6,0) -- (0.8,0);

        \node[
            below=2pt
        ] at (-0.9,0)
        {$ax+by+c'=0$};

        % Middle parallel line (equidistant)
        \draw[
            MPSecondary,
            densely dotted,
            line width=1pt
        ] (-2.6,0.575) -- (-0.1,0.575);

        \node[
            right=2pt
        ] at (-0.1,0.575)
        {$ax+by+\dfrac{c+c'}{2}=0$};

    \end{tikzpicture}
\end{center}

% ==========================================
% EXAMPLE
% ==========================================

\begin{example}
    معادلۀ یکی از محورهای تقارن مستطیلی که دو ضلع آن
    روی خطوط
    \(3x-y+7=0\) و \(6x-2y-6=0\)
    می‌باشد، کدام است؟

    \begin{tasks}(2)
        \task \(3x-y-2=0\)
        \task \(3x-y-1=0\)
        \task \(3x-y+1=0\)
        \task \(3x-y+2=0\)
    \end{tasks}
\end{example}

% ==========================================
% AREA OF A TRIANGLE
% ==========================================

\begin{remark}
    مساحت مثلثی که مختصات سه رأس آن معلوم باشد،
    به صورت زیر به دست می‌آید:
\end{remark}

\noindent
\begin{minipage}[c]{0.59\textwidth}
    \centering

    \resizebox{\linewidth}{!}{$
        S=\frac{1}{2}
        \left|
        x_{_A}(y_{_B}-y_{_C})
        +x_{_B}(y_{_C}-y_{_A})
        +x_{_C}(y_{_A}-y_{_B})
        \right|
    $}
\end{minipage}
\hfill
\begin{minipage}[c]{0.12\textwidth}
    \centering

    \begin{tikzpicture}[
        >=Stealth,
        line cap=round
    ]

        \node (A) at (0,0.95) {$A$};
        \node (B) at (-0.65,-0.15) {$B$};
        \node (C) at (0.65,-0.15) {$C$};

        \draw[
            ->,
            MPPrimary,
            line width=0.85pt,
            bend right=22
        ] (A) to (B);

        \draw[
            ->,
            MPPrimary,
            line width=0.85pt,
            bend right=22
        ] (B) to (C);

        \draw[
            ->,
            MPPrimary,
            line width=0.85pt,
            bend right=22
        ] (C) to (A);

    \end{tikzpicture}
\end{minipage}
\hfill
\begin{minipage}[c]{0.22\textwidth}
    \centering

    \begin{tikzpicture}[
        line cap=round,
        line join=round
    ]

        % Triangle vertices
        \coordinate (A) at (1.1,2.1);
        \coordinate (B) at (0,0);
        \coordinate (C) at (2.3,0.35);

        % Triangle
        \draw[
            MPPrimary,
            line width=1.1pt
        ] (A) -- (B) -- (C) -- cycle;

        % Vertex labels
        \node[above=2pt] at (A) {$A$};
        \node[below left=1pt] at (B) {$B$};
        \node[below right=1pt] at (C) {$C$};

    \end{tikzpicture}
\end{minipage}

% ==========================================
% EXAMPLE: AREA OF A TRIANGLE FROM MIDPOINTS
% ==========================================

\vspace{0.5\baselineskip}

\noindent
\begin{minipage}[c]{0.64\textwidth}
    \raggedleft

    \begin{example}
        اگر نقاط \(M(3,2)\)، \(N(6,2)\) و
        \(P(4,-3)\) نقاط میانی اضلاع یک مثلث
        (\(\overset{\triangle}{ABC}\)) باشند،
        مساحت این مثلث
        (\(\overset{\triangle}{ABC}\))
        برابر است با:

        \begin{tasks}(2)
            \task \(6.5\)
            \task \(7.5\)
            \task \(26\)
            \task \(30\)
        \end{tasks}
    \end{example}

\end{minipage}
\hfill
\begin{minipage}[c]{0.30\textwidth}
    \centering

    \begin{tikzpicture}[
        line cap=round,
        line join=round
    ]

        % Vertices of the original triangle
        \coordinate (A) at (1.2,2.8);
        \coordinate (B) at (0,0);
        \coordinate (C) at (2.8,0);

        % Midpoints of the three sides
        \coordinate (M) at ($(A)!0.5!(B)$);
        \coordinate (N) at ($(A)!0.5!(C)$);
        \coordinate (P) at ($(B)!0.5!(C)$);

        % Shaded medial triangle
        \fill[
            pattern=north east lines,
            pattern color=MPSecondary
        ] (M) -- (N) -- (P) -- cycle;

        % Outer triangle
        \draw[
            MPPrimary,
            line width=1.1pt
        ] (A) -- (B) -- (C) -- cycle;

        % Inner triangle
        \draw[
            MPPrimary,
            line width=1.1pt
        ] (M) -- (N) -- (P) -- cycle;

        % Vertex labels
        \node[above=2pt] at (A) {$A$};
        \node[below left=1pt] at (B) {$B$};
        \node[below right=1pt] at (C) {$C$};

    \end{tikzpicture}

\end{minipage}


% ==========================================
% REFLECTION AND SYMMETRY OF POINTS
% ==========================================

% Reflection of a point about another point
\begin{minipage}[c]{0.68\textwidth}
    \begin{remark}[توجه]
        قرینۀ نقطۀ
        $\Point{A}{x_1}{y_1}$
        نسبت به نقطۀ
        $\Point{O}{\alpha}{\beta}$
        برابر است با:
        \[
            \Point{A'}{2\alpha-x_1}{2\beta-y_1}
        \]
    \end{remark}

    \begin{remark}[نتیجه]
        قرینۀ نقطۀ
        $\Point{A}{x_1}{y_1}$
        نسبت به مبدأ مختصات برابر است با:
        \(
            \Point{A'}{-x_1}{-y_1}
        \)
    \end{remark}
\end{minipage}
\hfill
\begin{minipage}[c]{0.28\textwidth}
    \centering
    \begin{tikzpicture}[
        x=0.8cm,
        y=0.8cm,
        line cap=round
    ]
        % Points
        \coordinate (Ap) at (-1.7,-1.7);
        \coordinate (O) at (0,0);
        \coordinate (A) at (1.7,1.7);

        % Line segments
        \draw[MPPrimary, thick] (O)--(A);
        \draw[MPPrimary, thick, dotted] (Ap)--(O);

        % Points
        \fill[MPPrimary] (Ap) circle (2pt);
        \fill[MPPrimary] (O) circle (2pt);
        \fill[MPPrimary] (A) circle (2pt);

        % Labels
        \node[above right] at (A) {$A$};
        \node[below left] at (Ap) {$A'$};
        \node[below right] at (O) {$\Point{O}{\alpha}{\beta}$};
    \end{tikzpicture}
\end{minipage}

\begin{example}
    قرینۀ نقطۀ \(A(1,2)\) نسبت به نقطۀ \(M(-1,4)\) را بیابید.
\end{example}

% ==========================================
% REFLECTION ABOUT THE COORDINATE AXES
% ==========================================

\begin{remark}[توجه]
    اگر نقطۀ
    $\Point{A}{x_0}{y_0}$
    را در نظر بگیریم، آن‌گاه:

    \noindent
    \begin{minipage}[t]{0.65\textwidth}
        \vspace{0pt}
        \begin{enumerate}[label=\arabic*, leftmargin=*]
            \setlength{\itemsep}{0.25\baselineskip}

            \item تصویر \(A\) روی محور \(x\)ها برابر است با:
            $\Point{A'}{x_0}{0}$

            \item تصویر \(A\) روی محور \(y\)ها برابر است با:
            $\Point{A'}{0}{y_0}$

            \item قرینۀ \(A\) نسبت به محور \(x\)ها برابر است با:
            $\Point{A'}{x_0}{-y_0}$

            \item قرینۀ \(A\) نسبت به محور \(y\)ها برابر است با:
            $\Point{A'}{-x_0}{y_0}$
        \end{enumerate}
    \end{minipage}\hfill
    \begin{minipage}[t]{0.31\textwidth}
        \vspace{0pt}
        \centering
        \begin{tikzpicture}[
            x=0.75cm,
            y=0.75cm,
            >=stealth,
            line cap=round
        ]
            % Coordinate axes
            \draw[->, MPPrimary, thick]
                (-2.8,0)--(3.3,0)
                node[below right] {$x$};

            \draw[->, MPPrimary, thick]
                (0,-1.6)--(0,2.1)
                node[above left] {$y$};

            % Points
            \coordinate (A) at (2.3,1);
            \coordinate (Ap) at (-2,1);
            \coordinate (Ax) at (2.3,-1);

            % Dotted projections
            \draw[MPSecondary, dotted, thick] (Ap)--(A);
            \draw[MPSecondary, dotted, thick] (A)--(Ax);

            % Points
            \fill[MPPrimary] (A) circle (2pt);
            \fill[MPPrimary] (Ap) circle (2pt);
            \fill[MPPrimary] (Ax) circle (2pt);

            % Labels
            \node[above right] at (A) {$A$};
            \node[above] at (Ap) {$A'$};
        \end{tikzpicture}
    \end{minipage}

    \begin{enumerate}[label=\arabic*, leftmargin=*]
        \setcounter{enumi}{4}
        \setlength{\itemsep}{0.25\baselineskip}

        \item قرینۀ \(A\) نسبت به نیمساز ناحیۀ اول و سوم برابر است با:
        $\Point{A'}{y_0}{x_0}$

        \item قرینۀ \(A\) نسبت به نیمساز ناحیۀ دوم و چهارم
        \((y=-x)\) برابر است با:
        $\Point{A'}{-y_0}{-x_0}$
    \end{enumerate}
\end{remark}


% ==========================================
% EXAMPLE
% ==========================================

\begin{example}
    اگر \(A'\) قرینۀ \(A(-7,1)\) نسبت به محور \(x\)ها و
    \(A''\) قرینۀ \(A\) نسبت به نیمساز ناحیۀ اول و سوم باشند،
    اندازۀ \(A'A''\) چقدر است؟

    \begin{tasks}(4)
        \task \(12\)
        \task \(10\)
        \task \(6\sqrt{2}\)
        \task \(8\sqrt{2}\)
    \end{tasks}
\end{example}

\begin{example}
    اگر دو نقطۀ \(A(1-a,2b-1)\) و \(A'(a,b-3)\)
    قرینۀ یکدیگر نسبت به خط \(x+y=0\) باشند،
    \((a,b)\) کدام است؟

    \begin{tasks}(2)
        \task \((2,-1)\)
        \task \((-2,1)\)
        \task \((-1,1)\)
        \task \((1,-1)\)
    \end{tasks}
\end{example}

\begin{example}
    قرینۀ نقطۀ \(A(3,2)\) نسبت به خط \(y=x-3\) کدام است؟

    \begin{tasks}(2)
        \task \((4,3)\)
        \task \((5,0)\)
        \task \(\left(\frac{11}{2},\frac{7}{2}\right)\)
        \task \(\left(\frac{13}{3},\frac{9}{4}\right)\)
    \end{tasks}
\end{example}